% ============================================================================
%  Annexe A — Extraits de code
%  Règle : n'y placer que ce qui est cité dans le corps du texte.
%  Toute annexe non référencée depuis un chapitre doit être supprimée.
% ============================================================================
\chapter{Extraits de code}
\label{ann:extraits-code}

\section{Définition de l'environnement de compilation}
\label{ann:dockerfile}

\aecrire{Extrait du fichier de définition de l'image de compilation, réduit
aux parties significatives : image de base, installation de la chaîne
d'outils à versions épinglées, construction multi-étapes.}

\begin{lstlisting}[language=docker,caption={Extrait de la définition de
  l'image de compilation}]
# TODO : coller l'extrait.
\end{lstlisting}

\section{Chaîne de construction applicative}
\label{ann:workflow-app}

\begin{lstlisting}[language=yaml,caption={Extrait de la chaîne de construction
  applicative : déclencheurs et matrice}]
# TODO : coller l'extrait.
\end{lstlisting}

\section{Configuration de l'analyse statique}
\label{ann:coverity}

\begin{lstlisting}[language=yaml,caption={Configuration des phases de capture,
  d'analyse et de remontée}]
# TODO : coller l'extrait.
\end{lstlisting}

\section{Configuration de construction de la distribution}
\label{ann:kas}

\begin{lstlisting}[language=yaml,caption={Fichiers de configuration déclarative
  de la construction embarquée}]
# TODO : coller l'extrait.
\end{lstlisting}

\section{Description de l'infrastructure d'exécution}
\label{ann:terraform}

\begin{lstlisting}[language=hcl,caption={Extrait de la description
  d'infrastructure}]
# TODO : coller l'extrait, anonymise.
\end{lstlisting}
